\subsection{Markov Chain Monte Carlo}
% Metropolis-Hastings, Gibbs sampling, convergence diagnostics, mixing.

Monte Carlo integration in Method \ref{method:basic-monte-carlo-integration} uses independent draws from the target density. Markov chain Monte Carlo instead constructs a Markov chain $X_0,X_1,\ldots$ whose stationary density is the target $f$. When the chain mixes well, sample averages from the chain are used as approximations to expectations under $f$.

\begin{method} \label{method:metropolis-hastings} \hlt{Metropolis-Hastings}\\
Let $f$ be a target density available up to proportionality, and let $q(\cdot\mid x)$ be a tractable proposal density.\\
Choose $X_0$ arbitrarily; given $X_i=x$, define $X_{i+1}$ as follows.
\begin{enumerate}[label=(\roman*)]
\setlength{\itemsep}{0pt}
\item Draw $Y\sim q(\cdot\mid x)$.
\item Set $r(x,Y)=\min\left\{\frac{f(Y)q(x\mid Y)}{f(x)q(Y\mid x)},1\right\}$.
\item Set $X_{i+1}=Y$ with probability $r(x,Y)$, and set $X_{i+1}=x$ otherwise.
\end{enumerate}
\hlt{Key Insight}: Only ratios of $f$ enter $r$, so unknown normalising constants cancel.
\end{method}

\begin{proposition} \label{prop:mh-detailed-balance} \hlt{Detailed Balance for Metropolis-Hastings}\\
Let $\Phi(x,y)=f(x)q(y\mid x)$ and assume $\Phi(x,y)>0$ iff $\Phi(y,x)>0$.\\
For $x\neq y$, define $\alpha(x,y)=\min\left\{\frac{\Phi(y,x)}{\Phi(x,y)},1\right\}$ when $\Phi(x,y)>0$ and $\alpha(x,y)=0$ otherwise.\\
Set $p(x,y)=q(y\mid x)\alpha(x,y)$. Then $f(x)p(x,y)=f(y)p(y,x)$ for $x\neq y$.
\begin{proof}
If $\Phi(x,y)=0$, the support condition gives $\Phi(y,x)=0$, hence $p(x,y)=p(y,x)=0$.\\
If $\Phi(x,y)>0$, then $f(x)p(x,y)=\Phi(x,y)\min\left\{\frac{\Phi(y,x)}{\Phi(x,y)},1\right\}=\min\{\Phi(x,y),\Phi(y,x)\}$.\\
Interchanging $x,y$ gives the same expression for $f(y)p(y,x)$, QED.\\
\hlt{Structural Commentary}: The acceptance rule symmetrizes the weighted proposal flow.\\
\hlt{Key Insight}: Both directions reduce to the same minimum of the two proposal fluxes.
\end{proof}
\end{proposition}

\begin{proposition} \label{prop:detailed-balance-implies-stationarity} \hlt{Detailed Balance Implies Stationarity}\\
Let $\pi(dx)=f(x)\,dx$ and let $P$ be a Markov kernel on state space $\mathsf{X}$.\\
If $\pi(dx)P(x,dy)=\pi(dy)P(y,dx)$ as measures on $\mathsf{X}\times\mathsf{X}$, then $\pi P=\pi$.
\begin{proof}
For measurable $A\subseteq\mathsf{X}$, $\pi P(A)=\int P(x,A)\,\pi(dx)=\int \mathbf{1}_A(y)\,\pi(dx)P(x,dy)$.\\
By detailed balance, this equals $\int \mathbf{1}_A(y)\,\pi(dy)P(y,dx)=\int_A P(y,\mathsf{X})\,\pi(dy)=\pi(A)$, proving $\pi P=\pi$.\\
\hlt{Structural Commentary}: Detailed balance makes the joint law of consecutive states symmetric.\\
\hlt{Key Insight}: Marginalising the symmetric joint measure gives the same one-step density.
\end{proof}
\end{proposition}

\begin{method} \label{method:random-walk-metropolis} \hlt{Random-Walk Metropolis-Hastings}
\begin{enumerate}[label=(\roman*)]
\setlength{\itemsep}{0pt}
\item Use the proposal $Y=X_i+\varepsilon_i$, where $\varepsilon_i$ has a symmetric density.
\item Then $q(y\mid x)=q(x\mid y)$, so the acceptance probability is $r(x,y)=\min\left\{\frac{f(y)}{f(x)},1\right\}$.
\item For $\varepsilon_i\sim N(0,b^2)$, small $b$ gives slow local moves, large $b$ gives frequent rejection.\\
Tune $b$ for adequate mixing.
\end{enumerate}
\end{method}

\begin{example} \label{ex:cauchy-random-walk-mh} \hlt{Cauchy Target}\\
For Cauchy target $f(x)=\frac{1}{\pi(1+x^2)}$ and symmetric random-walk proposal, $r(x,y)=\min\left\{\frac{1+x^2}{1+y^2},1\right\}$.\\
The proposal scale changes mixing efficiency, not the stationary target.
\end{example}

\begin{method} \label{method:independence-metropolis} \hlt{Independence Metropolis-Hastings}
\begin{enumerate}[label=(\roman*)]
\setlength{\itemsep}{0pt}
\item Use a state-independent proposal $Y\sim g$, so $q(y\mid x)=g(y)$.
\item The acceptance probability is $r(x,y)=\min\{\frac{f(y)g(x)}{f(x)g(y)},1\}$.
\item Choose $g$ as a tractable approximation to the target with sufficient mass on important regions.
\end{enumerate}
\end{method}

\begin{method} \label{method:gibbs-sampling} \hlt{Gibbs Sampling}
\begin{enumerate}[label=(\roman*)]
\setlength{\itemsep}{0pt}
\item Let $(X,Y)$ have joint density $f_{X,Y}$ with tractable conditionals $f_{X\mid Y}$ and $f_{Y\mid X}$.
\item Starting from $(X_0,Y_0)$, iterate $X_{n+1}\sim f_{X\mid Y}(\cdot\mid Y_n)$ and $Y_{n+1}\sim f_{Y\mid X}(\cdot\mid X_{n+1})$.
\item In higher dimension, update one coordinate or block at a time from its full conditional.
\end{enumerate}
\end{method}

\begin{method} \label{method:metropolis-within-gibbs} \hlt{Metropolis within Gibbs}
\begin{enumerate}[label=(\roman*)]
\setlength{\itemsep}{0pt}
\item If $f_{X_j\mid X_{-j}}$ is not directly simulable, update $X_j$ by Metropolis-Hastings conditional on $X_{-j}$.
\item In the bivariate case, propose $Z\sim q(\cdot\mid X_n)$ and accept $X_{n+1}=Z$ with probability $\min\{\frac{f(Z,Y_n)q(X_n\mid Z)}{f(X_n,Y_n)q(Z\mid X_n)},1\}$.
\item Then update $Y$ analogously, conditional on the newly updated $X$.
\end{enumerate}
\end{method}
